\section{Opérateur constitutif et reconstruction des canaux du vide}
\label{iii:sec:constitutiva}

Les grandeurs constitutives s'obtiennent au moyen d'une opération sur le transport angulaire. L'information de départ comprend les canaux, leurs projecteurs et l'orientation qui les distingue. Cette organisation permet de suivre séparément la construction de l'opérateur, l'évaluation de ses fonctions spectrales et sa réalisation dimensionnelle ultérieure.

\subsection{Réponse positive sur deux feuillets}

Dans la chambre orientée $x>y>0$ de la section~\ref{iii:sec:hojas}, soient
$P_\pm=(I\pm\mathcal R)/2$ et
$T=q_+P_++q_-P_-$, avec $0<q_+<q_-<1$. Les incidences construites par le
calendrier fixent les secteurs $90$ et $120$. La réponse constitutive est
\begin{equation}
 \mathcal V=(I-T^{90})(I-T^{120})^{-1}.
 \label{iii:eq:vacuum-op}
\end{equation}
Cette définition reçoit le transport provenant du TPK et conserve les
attributions de feuillet. La règle de réponse et son domaine font explicitement partie
de la construction.

Le quotient provient de la comparaison des deux incidences du calendrier sur un même transport, et non de l'ajustement séparé de quatre grandeurs. Si $S=T^{30}$ et $\Sigma_j(S)=I+S+\cdots+S^{j-1}$, la factorisation de $I-S^j$ donne
\begin{equation}
 \mathcal V\,\Sigma_4(S)=\Sigma_3(S),\qquad
 \mathcal V=\Sigma_3(S)\Sigma_4(S)^{-1}.
 \label{iii:eq:completion-response}
\end{equation}
Les facteurs $\Sigma_3$ et $\Sigma_4$ expriment respectivement les
secteurs $90=3\cdot30$ et $120=4\cdot30$. Comme le spectre de $S$ est
contenu dans $(0,1)$, $\Sigma_4(S)$ est inversible : l’équation de
complétion détermine une unique réponse pour le transport fixé.
Cette unicité concerne la règle constitutive déclarée et ses
incidences. L’identification électromagnétique ajoute l’attribution
orientée des feuillets et les droites dimensionnelles précisées plus
bas.

\begin{proposition}[Décomposition et positivité]
La réponse est bien définie et
\begin{equation}
 \mathcal V=r_+P_++r_-P_-,\qquad
 r_\pm=\frac{1-q_\pm^{90}}{1-q_\pm^{120}}.
 \label{iii:eq:rchannels}
\end{equation}
Ses valeurs propres satisfont $3/4<r_-<r_+<1$ dans cette chambre.
\end{proposition}
\begin{proof}
Chaque coefficient $1-q_\pm^{120}$ est positif. L'inverse de $I-T^{120}$ est la somme des projecteurs divisés par ces coefficients, ce qui prouve \eqref{iii:eq:rchannels}. Avec $s=q^{30}$, on obtient
\begin{equation}
 f(s)=\frac{1+s+s^2}{1+s+s^2+s^3},\qquad
 f'(s)=-\frac{s^2(3+2s+s^2)}{(1+s+s^2+s^3)^2}<0.
 \label{iii:eq:monotone}
\end{equation}
Les limites en $0$ et en $1$ sont $1$ et $3/4$, respectivement. La monotonie
et l’ordre des canaux concluent la preuve.
\end{proof}

\subsection{Quatre coordonnées d’une réponse commune}

On définit les évaluations normalisées
\begin{equation}
 \widehat\varepsilon=r_+^2,\qquad
 \widehat\mu=r_-^2,\qquad
 \widehat Z=\frac{r_-}{r_+},\qquad
 \widehat c=\frac1{r_+r_-}.
 \label{iii:eq:four}
\end{equation}
Le premier couple conserve les réponses quadratiques assignées à chaque feuillet.
Le produit recueille le mode commun ; le quotient distingue leur orientation relative.
On vérifie exactement
\begin{equation}
 \widehat\mu\widehat\varepsilon=\widehat c^{-2},\qquad
 \frac{\widehat\mu}{\widehat\varepsilon}=\widehat Z^2,
 \qquad
 r_+=\frac1{\sqrt{\widehat Z\widehat c}},\quad
 r_-=\sqrt{\frac{\widehat Z}{\widehat c}}.
 \label{iii:eq:four-identities}
\end{equation}
Les racines positives fixent l'inversion. Sous la conjugaison des feuillets, lue par rapport aux marques $P_\pm$ fixes, $\widehat\varepsilon$ et $\widehat\mu$ sont échangés, $\widehat Z$ est inversé et $\widehat c$ est conservé. La réponse est alors transportée vers la chambre $x>-y>0$ : les formules restent valables, avec les ordres $q_-<q_+$ et $r_+<r_-$ inversés. Par conséquent, le mode commun et l'orientation sont des données récupérables au moyen d'observables complémentaires.

\begin{proposition}[Détermination des évaluations quadratiques]
Soit $W=r_+P_++r_-P_-$ une réponse positive sur les deux feuillets.
Étant fixées les lectures de produit et de quotient
\[
 \widehat c^{-1}=\det W=r_+r_-,
 \qquad \widehat Z=r_-/r_+,
\]
il existe une unique paire positive $(\widehat\varepsilon,\widehat\mu)$
qui vérifie
$\widehat\mu\widehat\varepsilon=\widehat c^{-2}$ et
$\widehat\mu/\widehat\varepsilon=\widehat Z^2$.
C’est précisément $(r_+^2,r_-^2)$.
\end{proposition}
\begin{proof}
La positivité permet de prendre les racines sans ambiguïté :
\[
 \widehat\mu=\frac{\widehat Z}{\widehat c}
 =\frac{r_-}{r_+}r_+r_-=r_-^2,\qquad
 \widehat\varepsilon=\frac1{\widehat Z\widehat c}
 =\frac{r_+}{r_-}r_+r_-=r_+^2.
\]
Les expressions vérifient les deux relations et en sont les seules racines
positives.
\end{proof}

Ainsi, les carrés ne constituent pas deux choix indépendants ajoutés
aux canaux : ils sont déterminés par les lectures multiplicative et
orientée, une fois fixées les relations constitutives. Cette proposition
établit l’unicité au sein de ce système de règles. L’interprétation
des règles comme réponse électromagnétique conserve son contenu
supplémentaire de représentation physique.

La lecture déterminantale du transport et la réponse constitutive
doivent également conserver leur ordre. Pour $T=e^{-xI-y\mathcal R}$, on a
$\det T=e^{-2x}$, tandis que
\[
 \det\mathcal V=f(q_+^{30})f(q_-^{30})=\widehat c^{-1}.
\]
Les deux proviennent du même opérateur à deux feuillets, mais évaluent des fonctions différentes de son spectre. En général, appliquer d'abord le déterminant puis une fonction rationnelle n'équivaut pas à prendre le déterminant de cette fonction de l'opérateur. Maintenir les deux valeurs propres jusqu'à l'achèvement de la réponse conserve l'information d'orientation qu'un déterminant isolé ne distingue plus.

\subsection{Inversion cubique et récupération angulaire}

\begin{theorem}[Inversion constitutive dans le domaine contractif]
Pour chaque $r\in(3/4,1)$, il existe une unique $s\in(0,1)$ telle que $f(s)=r$.
Cette $s$ est l’unique racine dans cet intervalle de
\begin{equation}
 r s^3+(r-1)(1+s+s^2)=0.
 \label{iii:eq:cubic}
\end{equation}
La réponse étiquetée reconstruit $s_\pm$, $q_\pm=s_\pm^{1/30}$ et
\begin{equation}
 x=-\frac12\log(q_+q_-),\qquad
 y=\frac12\log\frac{q_-}{q_+}.
 \label{iii:eq:recover-angular}
\end{equation}
\end{theorem}
\begin{proof}
La dérivée et les limites de \eqref{iii:eq:monotone} prouvent que $f$ est une
bijection décroissante entre les intervalles indiqués. Multiplier $f(s)=r$
par son dénominateur produit \eqref{iii:eq:cubic}. La racine positive d’ordre
$30$ retrouve le canal. Enfin, la somme et la différence de
$\log q_\pm=-x\mp y$ donnent \eqref{iii:eq:recover-angular}.
\end{proof}

La reconstruction conserve les étiquettes de feuillet. Avec l’opérateur complet
$\mathcal V$, on conserve aussi les projecteurs et on applique la fonction
inverse par calcul spectral. Récupérer deux coordonnées angulaires est une
opération définie sur cette représentation ; la trace historique complète
du TPK conserve des données supplémentaires.

L’intervalle $r\in(3/4,1)$ et la fonction inverse de ce théorème appartiennent
à la normalisation interne de \eqref{iii:eq:four}. Si l’on dispose de
coordonnées exprimées dans un autre système de coordonnées d’unités, on retrouve d’abord
cette normalisation au moyen de la transformation de
l’équation~\eqref{iii:eq:dimensions} ; on applique ensuite l’inversion cubique.

\subsection{Coordonnées logarithmiques et raffinement}

Soient $\mathcal E=\log(r_+r_-)$ et $\mathcal B=\log(r_-/r_+)$. Alors
\begin{equation}
 \log\mathcal V=\tfrac12\mathcal E I-\tfrac12\mathcal B\mathcal R.
 \label{iii:eq:logvac}
\end{equation}
Cette décomposition exhibe la composante invariante et la composante orientée.
Dans l’amplification $\mathcal V_m=\mathcal V\otimes I_{9^m}$, les
espérances partielles normalisées vérifient
$E_m(\mathcal V_{m+1})=\mathcal V_m$. Pour tout polynôme $p$,
$p(\mathcal V_m)=p(\mathcal V)\otimes I_{9^m}$ ; d’où
$E_m(p(\mathcal V_{m+1}))=p(\mathcal V_m)$. La continuité uniforme
étend l’identité aux fonctions continues du spectre. Les projecteurs,
les puissances et le logarithme sont ainsi conservés dans la tour indiquée.

\subsection{Types dimensionnels des quatre évaluations}

Avec les droites d’action, de charge et de vitesse, munies de sections
$u_S,u_Q,u_C$, la réalisation dimensionnelle s’écrit
\begin{align}
 Z_{\rm vac}&=\widehat Z u_Su_Q^{-2},&
 c_{\rm vac}&=\widehat c u_C,\\
 \mu_{\rm vac}&=\widehat\mu u_Su_C^{-1}u_Q^{-2},&
 \varepsilon_{\rm vac}&=\widehat\varepsilon u_S^{-1}u_C^{-1}u_Q^2.
 \label{iii:eq:dimensions}
\end{align}
En multipliant et en divisant ces expressions, on préserve
$\mu_{\rm vac}\varepsilon_{\rm vac}=c_{\rm vac}^{-2}$ et
$\mu_{\rm vac}/\varepsilon_{\rm vac}=Z_{\rm vac}^2$.
La coordinatisation SI évalue ensuite les sections. La comparaison physique nécessite
d’exposer leur construction et leur normalisation avec l’évaluation ; les quatre
coordonnées normalisées de \eqref{iii:eq:four} conservent leur type adimensionnel.
